%This file makes an abstract for a Ph.D. thesis.
\begin{center}
ABSTRACT
\end{center}
\vspace{0.3in}
\begin{center}
THE $x$-DEPENDENCE OF THE SIVERS TMD FROM LATTICE QCD
\\
BY
\\
HARIPRASHAD RAVIKUMAR, B.Sc., M.Sc., M.S.
\end{center}
\vspace{0.3in}
\begin{center}
Doctor of Philosophy

New Mexico State University

Las Cruces, New Mexico, 2026 

Dr. Michael Engelhardt, Chairperson 
\end{center}
\vspace{0.3in}
%The next line produces an indented paragraph to start the document
 %unit.  The LaTeX defaults start most units without indentations.
\hspace{\parindent}
We compute the dependence of the Sivers shift on the quark longitudinal momentum fraction $x$ in lattice QCD with 2+1 flavors of clover-improved Wilson fermions at a lattice spacing of \mbox{$a\approx 0.114$~fm}. The transverse momentum dependent parton distributions (TMDs) are defined through quark bilinears with staple-shaped Wilson lines, and the generalized Sivers shift is a ratio of Lorentz-invariant amplitudes. Earlier lattice calculations were restricted to $b\cdot P=0$ and gave only the $x$-integrated shift. We scan the full accessible range of $b\cdot P$, the variable Fourier conjugate to $x$. This requires staples along off-axis lattice directions, which we obtain by interpolation. Using the symbolic-regression tool PySR, we systematically explored several ans\"atze that satisfy physical constraints, and we describe the amplitudes at four nucleon momenta with a global fit in $b\cdot P$, $b^2$, and the staple length $\eta$. After extrapolation to the light-cone limit $\hat\zeta\to\infty$, the Sivers shift $\langle k_y\rangle_{TU}$ at a transverse quark separation of $0.34$~fm peaks at about $0.21$~GeV at small $x$ and falls to near zero by $x\approx 0.5$, while $x\langle k_y\rangle_{TU}$ peaks at about $0.02$~GeV near $x\approx 0.15$.
%This is the end of the abstract.
